\section{Verification, reproducibility, and remaining scope}
\label{sec:verification}

\subsection{What was checked}

The exact reference code was run with Python 3 using only standard-library
integer and rational arithmetic. The results recorded in
\texttt{code/verification.json} include:
\begin{center}
\begin{tabular}{@{}p{0.68\linewidth}r@{}}
\toprule
Check & Cases\\
\midrule
Bounded binary square root versus an independent implementation & 10,385\\
Family $X^{2m}+X+1$, $m=2,\ldots,12$ & 11\\
Quartics $X^4+aX^2+b$, $a,b\in[-8,8]$, $\Delta\ne0$ & 284\\
Additional square-completion inputs & 2\\
Malformed or out-of-scope inputs rejected & 9\\
\bottomrule
\end{tabular}
\end{center}

For quartics, the certificate enumerator, signed-divisor enumerator,
and an independent direct-power search agree. The independent search
uses a larger horizontal window and Python's separate integer-square-root
routine. These are implementation checks. They are not the proof
that the window contains every point; that fact comes from the
universal arithmetic bound. The family theorem for every $m\ge2$ is
likewise proved symbolically, not inferred from its eleven tested values.

The binary square-root implementation has explicit fuel. For $n\ge0$
it starts at $\ell=0,h=n+1$, maintains
$\ell^2\le n<h^2$, and replaces an endpoint by the integer midpoint
according to a square comparison. The interval length is at most
halved up to rounding. For $n\ge1$ with bitlength $b$, initially
$h-\ell=n+1\le2^b$, so $b+1$ iterations suffice to reach adjacent
endpoints. The case $n=0$ begins with adjacent endpoints. This proves
termination with a primitive recursive bound; it is not an
unbounded perfect-square search.

\begin{samepage}
The Lean build used version 4.32.0, commit
\begin{center}\small\ttfamily
8c9756b28d64dab099da31a4c09229a9e6a2ef35.
\end{center}
\end{samepage}
Its reproducible output is in \texttt{verification/lean\_build.txt}.
The exact theorem names and axiom dependencies are printed there.
A compiled definition of an unproved target is reported as a
specification, not as a proof. No complete ProveIt repository build
was performed.

The mathematical arguments were separately reviewed for their
coefficient estimates, nonzero remainder requirement, square
completion, sign cases, geometric interpretation, and logical
witness claims. The coding discussion was clarified to require
canonical point codes. The symbol $\Delta=a^2-4b$ is the auxiliary
quartic parameter; it is not the full polynomial discriminant.
The full quartic discriminant is $16b\Delta^2$.

\subsection{Archive contents}

\begin{center}
\begin{tabular}{@{}p{0.37\linewidth}p{0.54\linewidth}@{}}
\toprule
Path & Purpose\\
\midrule
\texttt{paper/siegel\_ipra.pdf} & Complete report\\
\texttt{paper/*.tex} & Master source and included sections\\
\texttt{lean/*.lean} & Checked native modules, specifications, and audit\\
\texttt{lean/OBJECT\_PLAN.txt} & Detailed object-iPRA compilation obligations\\
\texttt{lean/build.sh} & Reproducible Lean build\\
\texttt{code/runge.py} & Exact certificate constructor and solvers\\
\texttt{code/verify.py} & Independent finite regression checks\\
\texttt{code/build\_pdf.py} & Ordinary or private-format LaTeX build\\
\texttt{verification/} & Build and document-validation records\\
\texttt{provenance/} & Pinned source audit and research references\\
\texttt{README.txt} & Entry point, scope, and reproduction commands\\
\texttt{SHA256SUMS} & Checksums of the delivered files\\
\bottomrule
\end{tabular}
\end{center}

From the extracted root directory, run:
\begin{lstlisting}
python3 code/verify.py
bash lean/build.sh
python3 code/build_pdf.py
\end{lstlisting}
Use the version pinned in \texttt{lean/lean-toolchain}. If Lean is
available at a custom path, set \texttt{SIEGEL\_LEAN\_BIN} to that
executable before running the build script. The PDF script uses
ordinary \texttt{pdflatex} when configured and can prepare a private
format from installed TeX sources when the system format is missing.
It runs without shell escape.

Verify \texttt{SHA256SUMS} before rebuilding, since regenerated
logs and metadata may change. The archive contains source and
evidence; it does not bundle a Lean distribution or claim to have
recompiled the upstream foundation.

\subsection{The exact unresolved conclusion}

The fully developed arithmetic theorem proves effective finiteness
for the certified Runge family and supplies actual complete lists.
The archive also contains compiled native Lean proofs of important
lemmas and one global positive-genus example. The object-iPRA
derivations remain to be implemented using the explicit induction
invariants.

For arbitrary smooth positive-genus affine curves over number fields,
the requested general iPRA proof is not provided. A positive complete
list or height-bound proof would require the missing general effective
argument. A cardinality proof can pursue the different finite-tuple
route discussed above, whose approximation and geometry dependencies
have not been established in iPRA here. This report does not prove
either general target, and it does not claim either one impossible.

